%%%%%%%%%%%%%%%%%%%%%%%%%%%%%%%%%%%%%%%%%%%%%%%%%%%%%%%%%%%%%%%%%%%%%%%%%%%%
%% Author template for Management Science (mnsc) for articles with no e-companion (EC)
%% Mirko Janc, Ph.D., INFORMS, mirko.janc@informs.org
%% ver. 0.95, December 2010
%%%%%%%%%%%%%%%%%%%%%%%%%%%%%%%%%%%%%%%%%%%%%%%%%%%%%%%%%%%%%%%%%%%%%%%%%%%%
\documentclass[mnsc,blindrev]{informs3}
%\documentclass[mnsc,nonblindrev]{informs3} % current default for manuscript submission

\OneAndAHalfSpacedXI
%%\OneAndAHalfSpacedXII % Current default line spacing
%%\DoubleSpacedXII
%%\DoubleSpacedXI

% If hyperref is used, dvi-to-ps driver of choice must be declared as
%   an additional option to the \documentclass. For example
%\documentclass[dvips,mnsc]{informs3}      % if dvips is used
%\documentclass[dvipsone,mnsc]{informs3}   % if dvipsone is used, etc.

% Private macros here (check that there is no clash with the style)

%https://tex.stackexchange.com/questions/400692/different-height-of-tikz-subfigures-despite-same-axis-width-height

\usepackage{booktabs}
\usepackage{tabularx}
\usepackage{colortbl}
\usepackage{array}
\usepackage{multirow,calc,amsmath}
\newlength\mylenA
\newlength\mylenB
\newlength\mylenC
\newlength\mylenD
\usepackage{tikz}
\usepackage{tcolorbox}
\usepackage{etoolbox}
\AtBeginEnvironment{tcolorbox}{\normalsize}

\usepackage{adjustbox}[valign=t]

\tcbuselibrary{skins, breakable}

\makeatletter
\tcbset{
  tabular/.style={
    boxsep=\z@,top=\z@,bottom=\z@,leftupper=\z@,rightupper=\z@,
    toptitle=1mm,bottomtitle=1mm,boxrule=0.5mm,
    fontupper=\small,
    before upper={\arrayrulecolor{tcbcolframe}\def\arraystretch{1.1}%
      \tcb@hack@currenvir\tabular{#1}},
    after upper=\endtabular\arrayrulecolor{black}},
}
\makeatother

\usepackage{color, colortbl}
\usepackage{graphicx}
%\usepackage{subfigure}
\usepackage{caption}
\usepackage{subcaption}
%\usepackage{tcolorbox}
\usepackage{physics}
\usepackage{dsfont}
\usepackage{booktabs}
\usepackage{amsmath}
\usepackage{amsthm}
\usepackage{hyperref}
\usepackage{makecell}
\usepackage{tikz}
\usepackage{pgfplots}
\usepackage{tikzviolinplots}
\usepgfplotslibrary{external}
\tikzexternalize
\usepackage{minted}
\usemintedstyle{gruvbox-light}
\usepackage{scontents}
\usepackage{statmath}

 
\renewcommand\theadalign{bc}
\renewcommand\theadfont{\bfseries}
\renewcommand\theadgape{\Gape[4pt]}
\renewcommand\cellgape{\Gape[4pt]}

\newcommand{\gray}{\rowcolor[gray]{.90}}
% Natbib setup for author-year style
\usepackage{natbib}
 \bibpunct[, ]{(}{)}{,}{a}{}{,}%
 \def\bibfont{\small}%
 \def\bibsep{\smallskipamount}%
 \def\bibhang{24pt}%
 \def\newblock{\ }%
 \def\BIBand{and}%
\def\pipeline{prediction insight pipeline}
\newcommand{\tw}[1]{\textcolor{red}{Tong: #1}}

\def\checkmark{\tikz\fill[scale=0.4](0,.35) -- (.25,0) -- (1,.7) -- (.25,.15) -- cycle;} 

% \newtheorem{theorem}{Theorem}
% \newtheorem{corollary}{Corollary}[theorem]
% \newtheorem{lemma}[theorem]{Lemma}
% \newtheorem{remark}{Remark}
% \newtheorem{proof}{Proof}


%% Setup of theorem styles. Outcomment only one.
%% Preferred default is the first option.
\TheoremsNumberedThrough     % Preferred (Theorem 1, Lemma 1, Theorem 2)
%\TheoremsNumberedByChapter  % (Theorem 1.1, Lema 1.1, Theorem 1.2)
\ECRepeatTheorems

%% Setup of the equation numbering system. Outcomment only one.
%% Preferred default is the first option.
\EquationsNumberedThrough    % Default: (1), (2), ...
%\EquationsNumberedBySection % (1.1), (1.2), ...

% For new submissions, leave this number blank.
% For revisions, input the manuscript number assigned by the on-line
% system along with a suffix ".Rx" where x is the revision number.
\MANUSCRIPTNO{MS-0001-1922.65}

%%%%%%%%%%%%%%%%
\begin{document}
%%%%%%%%%%%%%%%%

% Outcomment only when entries are known. Otherwise leave as is and
%   default values will be used.
%\setcounter{page}{1}
%\VOLUME{00}%
%\NO{0}%
%\MONTH{Xxxxx}% (month or a similar seasonal id)
%\YEAR{0000}% e.g., 2005
%\FIRSTPAGE{000}%
%\LASTPAGE{000}%
%\SHORTYEAR{00}% shortened year (two-digit)
%\ISSUE{0000} %
%\LONGFIRSTPAGE{0001} %
%\DOI{10.1287/xxxx.0000.0000}%

% Author's names for the running heads
% Sample depending on the number of authors;
% \RUNAUTHOR{Jones}
% \RUNAUTHOR{Jones and Wilson}
% \RUNAUTHOR{Jones, Miller, and Wilson}
% \RUNAUTHOR{Jones et al.} % for four or more authors
% Enter authors following the given pattern:
%\RUNAUTHOR{}

% Title or shortened title suitable for running heads. Sample:
% \RUNTITLE{Bundling Information Goods of Decreasing Value}
% Enter the (shortened) title:
\RUNTITLE{The Illusion of Interpretation}

% Full title. Sample:
% \TITLE{Bundling Information Goods of Decreasing Value}
% Enter the full title:
%change title!!!
\TITLE{The Illusion of Interpretation: Post Hoc Explanations Aren't a Silver Bullet for Business Research}

% Block of authors and their affiliations starts here:
% NOTE: Authors with same affiliation, if the order of authors allows,
%   should be entered in ONE field, separated by a comma.
%   \EMAIL field can be repeated if more than one author
\ARTICLEAUTHORS{%
\AUTHOR{Snidely Slippery}
\AFF{Department of Bread Spread Engineering, Dairy University, Cowtown, IL 60208, \EMAIL{slippery@dairy.edu}} %, \URL{}}
\AUTHOR{Marg Arinella}
\AFF{Institute for Food Adulteration, University of Food Plains, Food Plains, MN 55599, \EMAIL{m.arinella@adult.ufp.edu}}
% Enter all authors
} % end of the block

% palette for colorblindness https://davidmathlogic.com/colorblind/#%23000000-%23E69F00-%2356B4E9-%23009E73-%23F0E442-%230072B2-%23D55E00-%23CC79A7






\bibliographystyle{informs2014}
\bibliography{refs}

\newpage

\end{document}
%%%%%%%%%%%%%%%%%

